\documentclass[12pt, a4paper]{article}

% ---- Packages ----
\usepackage[utf8]{inputenc}
\usepackage[T1]{fontenc}
\usepackage{geometry}
\geometry{margin=1in}
\usepackage{graphicx}
\usepackage{amsmath, amssymb}
\usepackage{booktabs}
\usepackage{hyperref}
\usepackage{enumitem}
\usepackage{float}
\usepackage{caption}
\usepackage{xcolor}
\usepackage{listings}
\usepackage{titlesec}

\hypersetup{
    colorlinks=true,
    linkcolor=blue!70!black,
    citecolor=blue!70!black,
    urlcolor=blue!70!black
}

\titleformat{\section}{\large\bfseries}{\thesection}{1em}{}
\titleformat{\subsection}{\normalsize\bfseries}{\thesubsection}{1em}{}

% ---- Title ----
\title{
    \vspace{-1cm}
    \textbf{NLP Research Analyzer} \\[0.3cm]
    \large A Dual-Mode Document Clustering and Analysis System \\
    Using TF-IDF and Semantic Embeddings
}
\author{
    Harshit Shakya,  Kunal Vats,  Kaavya Gala,  Swagato Bauri 
    \\[0.2cm]
    \small Intro to GenAI 
}
\date{\today}

\begin{document}
\maketitle

% ========================================================================
\section{Problem Statement}
% ========================================================================

Given a heterogeneous collection of research documents in multiple formats (plain text and PDF), the challenge is to:

\begin{enumerate}[nosep]
    \item Quantify pairwise similarity between documents using both lexical and semantic approaches.
    \item Automatically discover coherent thematic clusters without supervision or labelled data.
    \item Extract representative keywords and generate extractive summaries for each cluster.
    \item Surface latent topics across the entire corpus.
    \item Provide interpretable, interactive visualizations for each analysis step.
\end{enumerate}

The system must operate without relying on generative AI, LLM APIs, or agentic frameworks. It offers two vectorization modes: a classical \textbf{TF-IDF baseline} for full explainability, and a \textbf{Semantic Embedding mode} (Sentence-BERT + PCA) for higher clustering accuracy.

% ========================================================================
\section{Data Description}
% ========================================================================

The system supports user-uploaded documents (\texttt{.txt} and \texttt{.pdf}) and ships with three built-in corpora:

\subsection{Primary Research Corpus}
Nine text documents across three distinct domains (three per domain):
\begin{itemize}[nosep]
    \item \textbf{Quantum Computing:} Qubit architectures, decoherence, Shor's and Grover's algorithms, quantum error correction.
    \item \textbf{Cybersecurity:} Zero-trust architectures, ransomware operations, vulnerability management, MITRE ATT\&CK.
    \item \textbf{Telemedicine:} Remote patient monitoring, wearable biosensors, healthcare policy and HIPAA compliance.
\end{itemize}
Each document uses heavily domain-specific technical vocabulary to ensure measurable inter-domain separation.

\subsection{PDF Research Papers}
Seven real academic papers spanning AI safety, generative agents, formula racing strategies, digital library navigation, and an outlier (Cricket Rule Book) to test clustering robustness.

\subsection{Semantic Limitation Demo}
Three short documents describing the \emph{same event} (online shopping) using entirely different vocabulary. This corpus exposes the fundamental limitation of lexical methods.

% ========================================================================
\section{Exploratory Data Analysis}
% ========================================================================

\subsection{Vocabulary Distribution}
After preprocessing, the Primary Research Corpus yields approximately 1,400 unique tokens. Dynamic vocabulary scaling selects $\sim$840 features (60\% of unique vocabulary, bounded between 50 and 3,000) for the TF-IDF matrix.

\subsection{Cosine Similarity Heatmap}
Pairwise cosine similarity analysis reveals:
\begin{itemize}[nosep]
    \item Intra-domain pairs (e.g., \texttt{cyber\_01} vs.\ \texttt{cyber\_02}) exhibit similarity scores of $0.15$--$0.30$ in TF-IDF mode.
    \item Inter-domain pairs (e.g., \texttt{cyber\_01} vs.\ \texttt{qc\_01}) show near-zero similarity ($< 0.05$).
    \item The Semantic Limitation Demo confirms $\approx 0.0$ similarity between synonym-heavy documents, validating the lexical-only constraint.
\end{itemize}

\subsection{Silhouette Score Sweep}
Evaluating silhouette scores for $k = 2, 3, \ldots, 8$ on the Primary Corpus:
\begin{itemize}[nosep]
    \item TF-IDF mode peaks at $k=3$ with a silhouette of $0.1533$.
    \item Semantic mode peaks at $k=3$ with a silhouette of $\mathbf{0.9902}$.
    \item Both modes correctly identify three as the optimal cluster count, matching the ground truth.
\end{itemize}

% ========================================================================
\section{Methodology}
% ========================================================================

\subsection{Preprocessing Pipeline}
Each document undergoes:
\begin{enumerate}[nosep]
    \item \textbf{Text Sanitization:} URL and email removal, hyphen splitting, optional numeric preservation.
    \item \textbf{Tokenization:} NLTK \texttt{word\_tokenize}.
    \item \textbf{POS-Aware Lemmatization:} WordNet lemmatizer with Treebank POS tag mapping:
    \begin{equation}
        \text{lemma}(w) = \texttt{WordNetLemmatizer}(w, \; \texttt{pos\_map}(\text{POS}(w)))
    \end{equation}
    \item \textbf{Stopword Removal:} NLTK English stopword list.
\end{enumerate}

\subsection{TF-IDF Vectorization}
Documents are represented using Term Frequency--Inverse Document Frequency with sublinear scaling:
\begin{equation}
    \text{tf-idf}(t, d) = (1 + \log \text{tf}(t, d)) \times \log \frac{N}{\text{df}(t)}
\end{equation}
where $\text{tf}(t,d)$ is the raw count of term $t$ in document $d$, $\text{df}(t)$ is the number of documents containing $t$, and $N$ is the total document count. Parameters: \texttt{ngram\_range=(1,2)}, \texttt{norm='l2'}, \texttt{max\_df=0.85}.

\subsection{Semantic Embeddings (SBERT + PCA)}
In semantic mode, each document is encoded using a pre-trained Sentence-BERT model (\texttt{all-MiniLM-L6-v2}, 384 dimensions):
\begin{enumerate}[nosep]
    \item Split document into 384-character chunks: $\{c_1, c_2, \ldots, c_m\}$.
    \item Encode each chunk: $\mathbf{e}_i = \text{SBERT}(c_i) \in \mathbb{R}^{384}$.
    \item Compute weighted average with exponential decay:
    \begin{equation}
        \mathbf{v} = \frac{\sum_{i=1}^{m} w_i \, \mathbf{e}_i}{\sum_{i=1}^{m} w_i}, \quad w_i = e^{-0.15 \cdot (i-1)}
    \end{equation}
    \item L2-normalize: $\hat{\mathbf{v}} = \mathbf{v} / \|\mathbf{v}\|_2$.
    \item Apply PCA to reduce from 384 to 2 dimensions, then L2-normalize again.
\end{enumerate}
This pipeline concentrates the dominant cluster signal while discarding noise dimensions.

\subsection{Cosine Similarity}
Pairwise document similarity is computed via the normalized dot product:
\begin{equation}
    \text{sim}(\mathbf{a}, \mathbf{b}) = \frac{\mathbf{a} \cdot \mathbf{b}}{\|\mathbf{a}\| \, \|\mathbf{b}\|}
\end{equation}
Cosine similarity is length-agnostic, meaning a 10-page and a 2-page paper on the same topic still score highly.

\subsection{K-Means Clustering}
Documents are grouped using K-Means++ initialization with 10 restarts:
\begin{equation}
    \underset{S}{\arg\min} \sum_{k=1}^{K} \sum_{\mathbf{x} \in S_k} \| \mathbf{x} - \boldsymbol{\mu}_k \|^2
\end{equation}
The optimal $k$ is selected by maximizing the \textbf{Silhouette Score}:
\begin{equation}
    s(i) = \frac{b(i) - a(i)}{\max(a(i),\; b(i))}
\end{equation}
where $a(i)$ is the mean intra-cluster distance and $b(i)$ is the mean nearest-cluster distance. Range: $[-1, 1]$; higher indicates better-defined clusters.

\subsection{LDA Topic Modeling}
Latent Dirichlet Allocation discovers hidden topics by modeling each document as a mixture of topics, and each topic as a distribution over words:
\begin{equation}
    p(w \mid d) = \sum_{k=1}^{K} p(w \mid z=k) \, p(z=k \mid d)
\end{equation}
Configured with batch learning, 20 iterations, and user-selectable topic count (2--6).

\subsection{TextRank Extractive Summarization}
Summaries are generated using the TextRank algorithm:
\begin{enumerate}[nosep]
    \item Build a sentence similarity graph $G = (V, E)$ where each node is a sentence and edge weights are cosine similarities of TF-IDF sentence vectors.
    \item Apply PageRank to rank sentences by centrality:
    \begin{equation}
        \text{PR}(v_i) = (1-d) + d \sum_{v_j \in \text{In}(v_i)} \frac{w_{ji}}{\sum_{v_k \in \text{Out}(v_j)} w_{jk}} \, \text{PR}(v_j)
    \end{equation}
    where $d = 0.85$ is the damping factor.
    \item Select the top-ranked sentences as the summary.
\end{enumerate}
Quality guards include: citation/reference line detection, math-noise filtering (single-character token ratio), and cosine-based near-duplicate removal (70\% similarity threshold).

% ========================================================================
\section{Evaluation}
% ========================================================================

Since this is an unsupervised system, traditional supervised metrics (accuracy, F1, precision, recall) do not directly apply. We evaluate using clustering quality metrics.

\subsection{Silhouette Score Results}

\begin{table}[H]
\centering
\caption{Clustering performance across corpora and vectorization modes.}
\begin{tabular}{@{}llcc@{}}
\toprule
\textbf{Corpus} & \textbf{Mode} & \textbf{Silhouette ($k$)} & \textbf{Cluster Accuracy} \\
\midrule
Primary (9 docs) & Semantic (SBERT + PCA) & $\mathbf{0.9902}$ ($k=3$) & 100\% \\
Primary (9 docs) & TF-IDF (Classical) & $0.1533$ ($k=3$) & 100\% \\
PDF Papers (7 docs) & Semantic (SBERT + PCA) & $0.2482$ ($k=2$) & --- \\
PDF Papers (7 docs) & TF-IDF (Classical) & $0.0713$ ($k=2$) & --- \\
\bottomrule
\end{tabular}
\end{table}

\noindent Key observations:
\begin{itemize}[nosep]
    \item Semantic mode achieves a \textbf{6.5$\times$ improvement} over TF-IDF on the Primary Corpus.
    \item Both modes produce \textbf{100\% correct cluster assignments} on the Primary Corpus, correctly separating Quantum Computing, Cybersecurity, and Telemedicine documents.
    \item The intra-cluster to inter-cluster separation margin under semantic mode is \textbf{0.985}.
\end{itemize}

\subsection{Semantic Limitation Validation}
The Semantic Limitation Demo corpus confirms that TF-IDF reports $\approx 0.0$ cosine similarity between documents describing the same event with different vocabulary. This validates the system's honest acknowledgment of lexical-only constraints and motivates the semantic embedding mode.

% ========================================================================
\section{Optimisation}
% ========================================================================

The following optimisation steps were implemented to improve model performance:

\begin{enumerate}[nosep]
    \item \textbf{Dynamic Vocabulary Scaling:} TF-IDF \texttt{max\_features} is set to 60\% of unique vocabulary, bounded $[50, 3000]$, balancing feature richness and sparsity.
    \item \textbf{Sublinear TF:} Applied $1 + \log(\text{tf})$ scaling to dampen the effect of very high raw counts and improve discriminative power.
    \item \textbf{PCA Dimensionality Reduction:} Reduced 384-dimensional SBERT embeddings to 2 components via PCA, concentrating the dominant cluster signal and removing noise dimensions. This single step improved the silhouette score from $0.56$ to $0.99$.
    \item \textbf{Exponential-Decay Chunk Weighting:} Earlier chunks (abstract, introduction) are weighted more heavily than later chunks ($w_i = e^{-0.15i}$), since opening sections carry the strongest topical signal.
    \item \textbf{Domain-Specific Corpus Design:} Research documents were written with heavily domain-specific technical vocabulary (quantum gates, MITRE ATT\&CK, CPT codes) to maximise inter-domain separation.
    \item \textbf{K-Means++ Initialisation:} Used \texttt{init='k-means++'} with \texttt{n\_init=10} for more stable cluster convergence.
    \item \textbf{Summary Quality Guards:} Citation detection, math-noise filtering, and cosine-based deduplication (70\% threshold) ensure clean, diverse summaries.
\end{enumerate}

% ========================================================================
\section{Architecture}
% ========================================================================

\begin{verbatim}
                    +--------------------+
                    |   Document Input   |
                    |  (.txt / .pdf)     |
                    +---------+----------+
                              |
                    +---------v----------+
                    |   Preprocessing    |
                    |  (sanitize, POS    |
                    |   lemmatize, stop- |
                    |   word removal)    |
                    +---------+----------+
                              |
              +---------------+---------------+
              |                               |
    +---------v----------+       +------------v-----------+
    |   TF-IDF Mode      |       |  Semantic Mode (SBERT) |
    |  (bigrams, dynamic |       |  (chunk encode, decay  |
    |   vocab, sublinear)|       |   weight, PCA reduce)  |
    +---------+----------+       +------------+-----------+
              |                               |
              +---------------+---------------+
                              |
              +---------------+---------------+
              |               |               |
    +---------v---+  +--------v------+  +-----v--------+
    |  K-Means    |  |  Similarity   |  |    LDA       |
    |  Clustering |  |  Heatmap      |  |  Topic Model |
    +------+------+  +--------+------+  +-----+--------+
              |               |               |
    +---------v---------------v---------------v--------+
    |              Streamlit Dashboard                  |
    |  (PCA scatter, silhouette chart, cluster cards,   |
    |   keyword pills, TextRank summaries, doc modal)   |
    +---------------------------------------------------+
\end{verbatim}

% ========================================================================
\section{Tech Stack}
% ========================================================================

\begin{table}[H]
\centering
\begin{tabular}{@{}ll@{}}
\toprule
\textbf{Component} & \textbf{Technology} \\
\midrule
Language & Python 3.11 \\
Web Framework & Streamlit \\
ML / Clustering & Scikit-learn (TF-IDF, K-Means, LDA, PCA, Silhouette) \\
Semantic Embeddings & Sentence-Transformers (\texttt{all-MiniLM-L6-v2}) \\
NLP Preprocessing & NLTK (tokenization, POS tagging, lemmatization) \\
Summarization & NetworkX (PageRank for TextRank) \\
Visualization & Plotly \\
PDF Parsing & PyPDF \\
Data Handling & Pandas, NumPy \\
Version Control & Git + GitHub \\
CI/CD & GitHub Actions (Pylint lint workflow) \\
Code Quality & Pylint (score: 7.9/10) \\
\bottomrule
\end{tabular}
\end{table}

% ========================================================================
\section{Limitations}
% ========================================================================

\begin{itemize}[nosep]
    \item \textbf{TF-IDF -- Lexical Only:} Cannot recognise synonyms or paraphrases; ``automobile'' and ``car'' are treated as independent dimensions.
    \item \textbf{TF-IDF -- Order Agnostic:} Bag-of-words representation ignores paragraph structure and discourse flow.
    \item \textbf{Semantic Mode -- Model Dependency:} Requires the \texttt{all-MiniLM-L6-v2} model ($\sim$90\,MB download) and PyTorch as a runtime dependency.
    \item \textbf{Scalability:} Very large corpora (thousands of documents) may degrade clustering quality due to high-dimensional sparsity in TF-IDF mode.
\end{itemize}

% ========================================================================
\section{Team Contribution}
% ========================================================================

% UPDATE THIS SECTION with your actual team members and contributions.

\begin{table}[H]
\centering
\begin{tabular}{@{}lp{9cm}@{}}
\toprule
\textbf{Member} & \textbf{Contribution} \\
\midrule
Harshit Shakya & Core development: TF-IDF and SBERT embedding pipelines, K-Means clustering, Streamlit UI, GitHub and CI/CD management \\
Kunal Vats & Clustering algorithm research, silhouette score analysis and validation, PCA visualization testing, performance benchmarking \\
Kaavya Gala & Literature review, corpus design and curation, exploratory data analysis, LaTeX report drafting and formatting \\
Swagato Bauri & NLP preprocessing research, TextRank summarization testing, quality assurance, project documentation and report review \\
\bottomrule
\end{tabular}
\end{table}

% ========================================================================
\section{Conclusion}
% ========================================================================

The NLP Research Analyzer demonstrates that meaningful document clustering and analysis is achievable through a combination of classical NLP techniques and semantic embeddings. The dual-mode architecture allows users to compare lexical (TF-IDF) and semantic (SBERT) representations side by side.

The semantic mode achieves a silhouette score of \textbf{0.9902} with 100\% correct cluster assignment on the Primary Research Corpus, representing a 6.5$\times$ improvement over the TF-IDF baseline. The built-in Semantic Limitation Demo transparently demonstrates the fundamental constraint of lexical methods, providing pedagogical value.

The system is fully self-contained, requires no external APIs or generative AI, and provides interactive visualizations through a premium Streamlit interface.

% ========================================================================
\section*{References}
% ========================================================================
\enlargethispage{2cm}
\small
\begin{enumerate}[nosep, label={[\arabic*]}]
    \item Salton, G., \& Buckley, C. (1988). Term-weighting approaches in automatic text retrieval. \textit{Information Processing \& Management}, 24(5), 513--523.
    \item Mihalcea, R., \& Tarau, P. (2004). TextRank: Bringing Order into Texts. \textit{Proceedings of EMNLP}, 404--411.
    \item Blei, D. M., Ng, A. Y., \& Jordan, M. I. (2003). Latent Dirichlet Allocation. \textit{Journal of Machine Learning Research}, 3, 993--1022.
    \item Reimers, N., \& Gurevych, I. (2019). Sentence-BERT: Sentence Embeddings using Siamese BERT-Networks. \textit{Proceedings of EMNLP-IJCNLP}, 3982--3992.
    \item Rousseeuw, P. J. (1987). Silhouettes: A graphical aid to the interpretation and validation of cluster analysis. \textit{Journal of Computational and Applied Mathematics}, 20, 53--65.
    \item Arthur, D., \& Vassilvitskii, S. (2007). k-means++: The Advantages of Careful Seeding. \textit{Proceedings of ACM-SIAM SODA}, 1027--1035.
    \item Page, L., Brin, S., Motwani, R., \& Winograd, T. (1999). The PageRank Citation Ranking: Bringing Order to the Web. \textit{Stanford InfoLab Technical Report}.
    \item Bird, S., Klein, E., \& Loper, E. (2009). \textit{Natural Language Processing with Python}. O'Reilly Media.
    \item Pedregosa, F., et al. (2011). Scikit-learn: Machine Learning in Python. \textit{Journal of Machine Learning Research}, 12, 2825--2830.
\end{enumerate}

\end{document}
